\documentclass{article}
\usepackage{amsmath}
\usepackage{hyperref}
\usepackage[toc,section=section]{glossaries}

\makeglossaries

% Define the glossary entries
\newglossaryentry{ampere}{
    name=Ampere,
    description={The unit of electric current, denoted by \(I\), measuring the flow of electric charge per second. 1 Ampere (A) = 1 Coulomb/second}
}

\newglossaryentry{voltage}{
    name=Voltage,
    description={The electrical potential difference between two points in a circuit, measured in Volts (V). It drives the electric current through a circuit.}
}

\newglossaryentry{ohm}{
    name=Ohm,
    description={The unit of electrical resistance, denoted by \(\Omega\), which opposes the flow of electric current. Defined by Ohm's Law as \(R = V/I\)}
}

\newglossaryentry{resistor}{
    name=Resistor,
    description={A passive electrical component that limits the flow of current in a circuit. Its resistance is measured in ohms \((\Omega)\).}
}

\newglossaryentry{current}{
    name=Current,
    description={The flow of electric charge through a conductor, measured in Amperes (A). Denoted by \(I\). Follows the relationship \(I = V/R\) from Ohm's Law.}
}

\newglossaryentry{power}{
    name=Power,
    description={The rate at which electrical energy is consumed or converted, measured in Watts (W). Calculated as \(P = VI = I^2 R\).}
}

\newglossaryentry{energy}{
    name=Energy,
    description={The capacity to do work, measured in Joules (J). Electrical energy is calculated as \(E = P \times t\), where \(t\) is time.}
}

\newglossaryentry{polarity}{
    name=Polarity,
    description={The orientation of electrical potential in a circuit. Positive and negative terminals indicate the direction of current flow.}
}

\newglossaryentry{circuit}{
    name=Circuit,
    description={A closed loop that allows electric current to flow, typically including components like resistors, capacitors, and power sources.}
}

\newglossaryentry{switch}{
    name=Switch,
    description={A device that opens or closes a circuit, allowing or preventing current to flow.}
}

\newglossaryentry{battery}{
    name=Battery,
    description={A device that stores chemical energy and converts it into electrical energy to provide voltage for a circuit.}
}

\newglossaryentry{direction}{
    name=Direction,
    description={The flow of electric current, conventionally considered from the positive to the negative terminal of a battery.}
}

% Basic Rules for Solving Circuits
\newglossaryentry{ohmslaw}{
    name=Ohm's Law,
    description={The fundamental relationship between voltage, current, and resistance in a circuit: \(V = IR\), where \(V\) is voltage, \(I\) is current, and \(R\) is resistance.}
}

\newglossaryentry{kirchhoffsv}{
    name=Kirchhoff's Voltage Law (KVL),
    description={The sum of all voltages around a closed loop in a circuit is zero. Mathematically, \(\sum V = 0\). This law is used to analyze the voltage drops across components.}
}

\newglossaryentry{kirchhoffsc}{
    name=Kirchhoff's Current Law (KCL),
    description={The sum of currents entering a junction equals the sum of currents leaving the junction. Mathematically, \(\sum I = 0\).}
}

\begin{document}

\title{Electrical Circuits Glossary and Basic Rules}
\author{Your Name}
\date{\today}
\maketitle

\section*{Glossary of Terms}

\printglossaries

\section*{Basic Rules for Solving Circuits}

\subsection*{Ohm's Law}
Ohm's Law states that the voltage (\(V\)) across a resistor is proportional to the current (\(I\)) through the resistor, given by \(V = IR\). This is fundamental to solving circuits.

\subsection*{Kirchhoff's Voltage Law (KVL)}
Kirchhoff's Voltage Law states that the sum of all the voltages around any closed loop in a circuit must equal zero. For example, in a simple series circuit:
\[
V_{\text{battery}} = V_1 + V_2 + V_3
\]
where \(V_1\), \(V_2\), and \(V_3\) are the voltage drops across the resistors.

\subsection*{Kirchhoff's Current Law (KCL)}
Kirchhoff's Current Law states that the total current entering a junction must equal the total current leaving the junction:
\[
I_{\text{in}} = I_{\text{out}}
\]
This is useful for analyzing parallel circuits.

\subsection*{Power in a Circuit}
Power in an electrical circuit is given by \(P = VI\) and can also be expressed as:
\[
P = I^2 R \quad \text{or} \quad P = \frac{V^2}{R}
\]
Power represents the rate at which energy is consumed by a resistor or other component.

\section*{Solving Circuits with Batteries, Switches, and Resistors}

1. **Step 1**: Identify all resistors, power sources (batteries), and switches.
2. **Step 2**: Use Ohm’s Law \(V = IR\) to calculate unknown values (current, resistance, voltage).
3. **Step 3**: Apply Kirchhoff’s Voltage Law (KVL) to any closed loops to find unknown voltages.
4. **Step 4**: Apply Kirchhoff’s Current Law (KCL) to any junctions to find unknown currents.
5. **Step 5**: Calculate the power consumed by each resistor using \(P = I^2 R\).

\end{document}
